% ECONOMICS_PROBLEM: {"id": "Q57-292", "title": "Biodiversity and ecosystem-service valuation", "jel_code": "Q57", "source_id": "OP-0117"}
% !TeX program = xelatex
\documentclass[11pt,a4paper]{article}
\usepackage[margin=23mm]{geometry}
\usepackage{fontspec}
\setmainfont{Latin Modern Roman}
\usepackage{amsmath,amssymb,mathtools}
\usepackage{xcolor,enumitem,booktabs,longtable,array,xurl,hyperref}
\definecolor{muted}{HTML}{53636C}
\hypersetup{colorlinks=true,urlcolor=blue,linkcolor=blue}
\setlength{\parindent}{0pt}
\setlength{\parskip}{5pt}
\newcommand{\R}{\mathbb R}
\newcommand{\E}{\mathbb E}
\newcommand{\N}{\mathbb N}
\newcommand{\ind}{\mathbf 1}
\newcommand{\Var}{\operatorname{Var}}
\newcommand{\Cov}{\operatorname{Cov}}
\newcommand{\supp}{\operatorname{supp}}
\DeclareMathOperator*{\argmin}{arg\,min}
\DeclareMathOperator*{\argmax}{arg\,max}
\newcommand{\entrymeta}[3]{{\sffamily\small\color{muted}\textbf{\texttt{#1}}\quad|\quad #2\par #3\par}}
\newcommand{\Problemref}[1]{\texttt{#1}}
\newcommand{\JELref}[2]{\texttt{#2} (JEL \texttt{#1})}
\begin{document}
\section*{Biodiversity and ecosystem-service valuation}
\textbf{Problem ID:} \texttt{Q57-292}\quad \textbf{JEL:} \texttt{Q57}\par
% BEGIN ECONOMICS STATEMENT
\label{problem:OP-0117}
\label{audited:ECON-058}


\label{atlas:prob:C7}

\noindent\textbf{Original identifier:} C7 (atlas item 117). \textbf{Field:} Environmental and climate economics. \textbf{Classification:} editorial research family with established restricted solutions; current open status not certified.

\subsection*{Definitions and assumptions}

Let ecosystem state $s_t\in\mathbb R_+^K$ record explicitly defined stocks, such as habitat areas and species populations, each in its own units. Ecological transitions are $s_{t+1}=F(s_t,a_t,\xi_{t+1})$, with land-use action $a_t$ and shocks $\xi_t$. Ecosystem services $q_t=G(s_t)$ enter household preferences or production functions. A specified economic equilibrium determines consumption $c_{h,t}$, and social welfare is $W(s_0)=\mathbb E[\sum_{t=0}^H\beta_t\sum_h\omega_hu_h(c_{h,t},q_t)]$. Fix weights $\omega_h$, discounting $\beta_t$, horizon, ecological thresholds, and the currency/location of a compensating consumption transfer with marginal welfare value $\lambda_0>0$.

\subsection*{Formal research question}

Identify economically meaningful values for marginal ecosystem changes and nonmarginal irreversible losses, with quantified uncertainty. Where differentiable, the shadow price of stock component $k$ is
\[
p_k(s_0)=\lambda_0^{-1}\frac{\partial W(s_0)}{\partial s_{0k}},
\]
in real currency per unit of that component. For a specified irreversible transition from regime $A$ to $B$, define a compensating transfer $T$ by $W(B;T)=W(A;0)$ whenever such a finite transfer exists. Characterize identified sets for $p_k$ and $T$ over preference, ecological and economic models consistent with market behavior, stated-preference evidence and ecological data. Determine when thresholds and poor substitutability imply discontinuity, lack of finite compensation or wide irreducible bounds; construct policy rankings robust to that ambiguity.

\subsection*{Interpretation and scope}

A monetary total of annual service flows is not a stock value or a causal marginal price for biodiversity loss. Species counts, functional diversity and habitat integrity are distinct state variables; exchange rates among them require ecological and welfare assumptions. Nonuse values must appear in preferences rather than being inferred from absent market transactions. Summing overlapping services can count the same contribution multiple times. At an extinction threshold a derivative may not summarize the relevant discrete policy change. Money-metric compensation can fail when essential services cannot be substituted; then a feasibility or welfare comparison is appropriate without manufacturing a finite price. Natural-capital accounts require stock-flow consistency and declared accounting boundaries.

\subsection*{Source audit and status}

[S1] is the verified 1997 global ecosystem-service valuation article; [S3] is the official 2021 Dasgupta Review. \emph{Polasky et al.} lacks title and year and remains unresolved [S2]; a verified 2005 biodiversity handbook chapter is listed only as a candidate, not silently substituted. These sources motivate valuation and accounting but do not establish universally identifiable biodiversity prices. The nonlinear stock-and-compensation formulation is editorial, with no certified current open-status assertion.

\subsection*{References}

\begin{enumerate}[label={[\textnormal{S}\arabic*]},leftmargin=*]

\item Robert Costanza, Ralph d’Arge, Rudolf de Groot, Stephen Farber, Monica Grasso, Bruce Hannon, Karin Limburg, Shahid Naeem, Robert V. O’Neill, Jose Paruelo, Robert G. Raskin, Paul Sutton and Marjan van den Belt (1997). \emph{The value of the world’s ecosystem services and natural capital}. Nature 387, 253–260. \url{https://doi.org/10.1038/387253a0}. \textbf{Locator:} Publisher or institutional abstract and bibliographic record. \textbf{Support and limits:} Global service-flow valuation exercise; total flows are not a marginal welfare price for a species or irreversible change.

\item \textbf{Unresolved original citation:} Polasky et al. (undated). Multiple relevant works exist. Polasky, Costello and Solow (2005), The Economics of Biodiversity, is a verified candidate, not an identification of the unspecified original. Candidate/verification route: \url{https://www.sciencedirect.com/science/article/pii/S1574009905030299}.

\item Partha Dasgupta (2021). \emph{The Economics of Biodiversity: The Dasgupta Review}. HM Treasury, final report. \url{https://www.gov.uk/government/publications/final-report-the-economics-of-biodiversity-the-dasgupta-review}. \textbf{Locator:} Publisher or institutional abstract and bibliographic record. \textbf{Support and limits:} Natural-capital and biodiversity framework; it does not deliver universally identifiable shadow prices.

\end{enumerate}

\subsection*{Integrated refinement: ECON-058}
{\sffamily\small\color{muted}Economic value of biodiversity in forest policy\par}
This refinement adopts the additional source conventions
(page~\pageref{audited:conventions}), subject to its local restrictions.
\textbf{Forest biodiversity under risk: ex-ante versus state-specific compensation.}
At fixed reference prices and strictly increasing income utility $v_i(y,b)$,
define deterministic ex-ante compensation, when it exists, by
\[
\mathbb E[v_i(y_i-CV_i^{EA},b_1)]=v_i(y_i,b_0).
\]
Alternatively solve the state-specific equality for $CV_i(\omega)$ under
$b_1(\omega)$. These differ under risk aversion even for the same mean ecological
change. Identify or bound the separately declared aggregates
$\sum_iCV_i^{EA}$ or $\sum_i\mathbb E[CV_i(\omega)]$, with population weights,
income domain, existence/range and uniqueness conditions, and common horizon and
currency. Include induced production and livelihood responses. Exclude carbon
and timber benefits counted elsewhere; identify ecological effects separately
from valuation. Threshold losses may lack finite compensation.
\subsection*{Supplied source audit for ECON-058}
Local [S1], [S2], etc. in this audit and the references immediately below
belong to ECON-058, independently of the original entry's reference numbering.
[S1]'s conclusion explicitly identifies the difficulty of quantifying economic biodiversity gains, beyond ecological impacts. [S2] corroborates the topic. The risk-specific compensation definitions are editorial refinements.
\subsection*{References for integrated refinement ECON-058}
{\footnotesize\raggedright
\par\noindent\textbf{[S1]} Francisco Costa, Allan Hsiao, Heitor Pellegrina, and Eduardo Souza-Rodrigues (2025). \emph{Deforestation}. VoxDevLit 18(1), September 2025.\par \textit{Verified locator / source role:} Conclusion/evidence-gap chapter at the linked primary review URL, inspected for the agenda. See the audit conclusion for distinctions between direct agenda support and editorial extensions.\par \url{https://www.voxdev.org/voxdevlit/deforestation/conclusion-and-evidence-gaps}\par\smallskip
\par\noindent\textbf{[S2]} Oliver Hanney / VoxDev (living compilation). \emph{Evidence gaps: Key unanswered questions in development economics}. Accessed 8 October 2026. The page currently covers 20 topics; the ``16-topics URL is a historical slug.\par \textit{Verified locator / source role:} Topic heading: Deforestation. Supports the broad research agenda; this is not a verbatim quotation or an attribution of the displayed mathematical target.\par \url{https://voxdev.org/topic/evidence-gaps-key-unanswered-questions-across-16-topics-development-economics}\par\smallskip
}

\subsection*{Common conventions: Source mathematical conventions}

Unless an entry states otherwise, agent and firm sets are finite. State and action spaces are measurable, strategies and outcome maps are measurable, probability laws are specified, and expectations exist for the targets under discussion. Infinite-horizon discount factors lie in $(0,1)$ and discounted objectives are finite. Norms, tolerances and complexity input sizes must be specified for computational comparisons. Constants or primitives that appear locally are inputs, not empirical facts asserted by this catalogue.

\textbf{Identification.} A statistical model $m\in\mathcal M$ implies an observable law $P_m$ and target $\theta(m)$. At law $P$, its identified set is
\[
\Theta(P;\mathcal M)=\{\theta(m):m\in\mathcal M,\ P_m=P\}.
\]
Point identification holds when this set is a singleton, or equivalently when $P_{m_1}=P_{m_2}$ implies $\theta(m_1)=\theta(m_2)$ for all compatible models. Sharp bounds mean bounds attained, or approached when the boundary is not attained, by compatible models. Writing this set is a definition, not a solution: the substantive task is to establish informative restrictions and derive or estimate the set. An unrestricted-confounding model may give only trivial bounds.

\textbf{Inference.} Identification, estimability and valid uncertainty quantification are separate questions. Fix $\alpha\in(0,1)$. A confidence procedure $C_n$ over a declared law class $\mathcal P$ has asymptotic uniform coverage at level $1-\alpha$ when
\[
\liminf_{n\to\infty}\inf_{P\in\mathcal P}\Pr_P\{\theta(P)\in C_n\}\geq1-\alpha.
\]
For set-identified targets the coverage event must be defined separately, for example coverage of the full identified set. If an entry uses identification-indexed models instead of uniquely indexed laws, the same coverage convention is applied over those models. Pointwise limits do not establish uniform validity, and asymptotic validity does not guarantee finite-sample precision. Sample sizes, dependence structures and weak-identification sequences are part of each application.

\textbf{Counterfactuals.} $Y_i(a)$ is a potential outcome under a specified intervention, including its timing, implementation and versions. It is not $Y_i$ conditional on voluntarily choosing $a$. A full intervention vector permits interference; shorthand scalar treatment notation does not silently impose no interference. Assignment, selection, attrition, anticipation and measurement must be specified. A claim about an instrument's local effect is not automatically a population policy effect.

\textbf{Economic equilibrium.} A counterfactual model includes best responses, constraints, feasibility, market clearing, state transitions and expectations consistent with the stated information. When several equilibria exist, either a declared selector is an input or the target is set-valued. Comparisons across policy regimes must use the stated selection convention. Reduced-form relationships are not presumed invariant after an equilibrium-changing intervention.

\textbf{Optimization and welfare.} Optimization is over the stated feasible set. If attainment is not established, $\sup$ or $\inf$ and $\varepsilon$-optimal choices replace an asserted maximizer or minimizer. For example, with value $V=\sup_{a\in\mathcal A}W(a)<\infty$, an $\varepsilon$-optimal set is $\{a\in\mathcal A:W(a)\geq V-\varepsilon\}$ for $\varepsilon>0$. Benefits, costs, risk and health or environmental outcomes can be aggregated only using declared units and conversion weights. Interpersonal utility weights, the treatment of fiscal transfers and foreign profits, discounting, and the behavioral welfare criterion are explicit inputs. A comparative-static derivative additionally requires existence and appropriate differentiability of the selected counterfactual.

\textbf{Sources.} Local labels [S1], [S2], and so forth restart in every question. Locators identify the relevant abstract, model, section or result at the level actually checked; a bibliographic or abstract-level verification is not represented as inspection of an entire proof. Working papers are distinguished from later publications and changing versions. Source summaries are paraphrased. All URL checks and editorial formulations refer to the audit date stated above.

\subsection*{Common conventions: Additional-source status and mathematical conventions}

\label{audited:conventions}

Of the 100 supplied ECON entries, 65 distinct problems appear here; 32 close
matches refine earlier entries, and three duplicate questions map to existing
entries. Appendix~\ref{app:audited-crosswalk} lists every destination.
The attachment's P/R/S/U distinctions and source audits are retained. Its
precise mathematical formalizations are editorial unless the individual audit
says otherwise. Candidate subsidy targets ECON-004--006 retain their U flags;
the resolved approval-core question ECON-007 updates OP-0202 above.
The following supplied conventions govern both this part and the attributed
ECON refinements elsewhere in the catalogue, subject to local restrictions.

\textbf{Parameterized questions.} A problem family lists inputs that must be fixed before an instance is evaluated: population, horizon, policies, information, data law, model class and welfare weights. These are required choices, not quantities the citations are assumed to specify. Undefined applications should not be treated as identified estimands.

\textbf{Indices and integrability.} Sets of agents, firms and regions are finite unless a population measure is explicitly used. \(i,j\) index units, \(r\) regions, \(t\) dates and \(h\) horizons. \(\mathbb E\) and \(\Pr\) refer to the declared population and law; \(\mathbf1\{A\}\) is an indicator. Every displayed expectation and discounted sum needs existence in the stated domain. Infinite horizons use a discount in \((0,1)\) and a convergence condition; finite horizons may use discount one.

\textbf{Optimization and existence.} A displayed \(\max\), \(\min\), or \(\arg\max\) is conditional on attainment. For finite feasible menus this is immediate; general spaces need continuity and compactness/coercivity or another existence argument. Without attainment, the target is the supremum/infimum and arbitrarily near-optimal feasible choices. \(\inf\varnothing=+\infty\). Empty feasibility and an unidentified target are different issues.

\textbf{Potential outcomes.} \(Y_i(z)\) is the outcome under a specified intervention, not the observed conditional mean among people choosing \(z\). In binary comparisons \(\Delta Y_i=Y_i(1)-Y_i(0)\). Specify assignment, treatment versions, compliance, information and observation horizon. Interference is allowed under a full policy vector; compressed exposure notation needs a justified mapping. No interference, consistency and overlap are substantive assumptions, not automatic consequences of notation.

\textbf{Populations and observation.} Fix baseline eligibility and population weights. Conditioning on post-policy employment, survival, relocation or program uptake can change the population being compared. Death is not ordinary loss to follow-up; outcomes truncated by death need a composite convention, joint survival target, or explicitly defined survivor stratum. Fertility-changing interventions need a population functional rather than an unexplained fixed descendant cohort. Measurement and missingness models must be supplied.

\textbf{Identification.} A structural model \(m\in\mathcal M\) implies observable law \(P_m\) and target \(\theta(m)\). Identification means
\[
 P_{m_1}=P_{m_2}\ \Longrightarrow\ \theta(m_1)=\theta(m_2)
 \quad\text{for all }m_1,m_2\in\mathcal M .
\]
Without identification, the target is
\[
 \Theta(P;\mathcal M)=\{\theta(m):m\in\mathcal M,\ P_m=P\}.
\]
Writing \(\theta(P)\) before establishing identification can hide incompatible latent models. Confidence sets additionally need a sampling law, confidence level, asymptotic sequence and uniformity class. Point identification, coverage and informative precision are separate properties.

\textbf{Mediation.} Total effects of randomized assignment are distinct from cross-world natural indirect effects. Belief/effort-channel effects require a meaningful mediator intervention and additional measurement, exclusion or mediation assumptions. Randomization of the initial policy alone does not supply those assumptions. Report bounds or interventional alternatives when the stronger target is unsupported.

\textbf{Equilibrium.} A counterfactual \(W(z)\), \(p(z)\), or \(\mu_t^z\) requires individual optimization, feasibility, government/resource budgets, market clearing and state-transition laws. State equilibrium existence and selection, or report ranges over compatible equilibria. Initial-state integration includes subsequent endogenous transitions; current-state integration must use the policy-dependent distribution. Reduced-form invariance after policy is not assumed.

\textbf{Welfare accounts.} Scores, utility, health, dollars and ecological quantities require explicit conversion before addition. Fees, taxes, subsidies, wages and extortion payments are transfers between parties; distributional weights can make them welfare relevant, but their face value is not automatically a resource loss. State ownership, geographic boundaries, foreign-income weights and fiscal finance. Do not add profits to household welfare already including dividends, or count land capitalization and the underlying benefit twice. A benefit-cost expression must distinguish gross benefits from net welfare and subtract every real cost exactly once.

\textbf{Derivatives and risk.} Log derivatives require positive arguments; local responses require admissible perturbations and specified equilibrium branches. Differentiation under expectations needs regularity/domination, and boundaries may allow only one-sided derivatives. Risk uses a specified law; ambiguity uses a declared nonempty law set on a common history space. Adaptive policies must be nonanticipative. A robust criterion is an editorial choice unless attributed explicitly.

\textbf{Scope overlaps.} Allocation, collective choice, contracts, household bargaining and coordinated policy overlap economics with optimization and game theory. They are retained because they appeared in the original catalogue; no claim of a strictly game-theory-free selection is made.
% END ECONOMICS STATEMENT
\end{document}
